
\section{ASP Symbolic Refinement Layer}
This appendix documents the symbolic reasoning layer used in the proposed neuro-symbolic anomaly-analysis framework. The symbolic component is implemented in Answer Set Programming (ASP) and is responsible for refining the raw anomaly candidates produced by the neural residual-based detector. Rather than generating anomalies independently, the ASP layer receives candidate anomalies and contextual facts from the Python pipeline and applies a compact set of domain rules to determine whether each anomaly should be retained or rejected.

The role of this appendix is to make the symbolic component of the framework fully transparent and reproducible. In particular, it documents the complete ASP rule base used by the system, explains the predicate structure that connects the neural and symbolic stages, and provides the exact rule file as implemented in the repository.

\section{ASP Predicate Structure and Integration Logic}

The symbolic refinement stage operates on facts generated from the anomaly table and the engineered hourly dataset. These facts are constructed in the Python integration script \texttt{src/07\_apply\_asp.py}. The script translates raw anomaly detections, time information, holiday indicators, and selected contextual variables into symbolic atoms that can be processed by the ASP solver.

The main predicate families used by the ASP layer are as follows:
\begin{itemize}
	\item \texttt{hour(T)}: declares a discrete hourly time point;
	\item \texttt{anomaly(G,T)}: indicates that signal \texttt{G} was flagged as anomalous at time \texttt{T};
	\item \texttt{holiday(T)}: indicates that time \texttt{T} falls on a public holiday;
	\item \texttt{pred(shortwave\_radiation,R,T)}: stores rounded shortwave-radiation information;
	\item \texttt{pred(wind\_mw,V,T)}: stores rounded wind-generation information;
	\item \texttt{pred(load\_mw,V,T)}: stores rounded actual-load information;
	\item \texttt{pred(load\_ref,R,T)}: stores the reference load used in contextual reasoning;
	\item \texttt{valid\_anomaly(G,T)}: marks a candidate anomaly retained after symbolic reasoning;
	\item \texttt{vetoed(G,T)}: marks a candidate anomaly rejected by the symbolic layer;
	\item \texttt{veto\_reason(T,R)}: records the symbolic reason associated with a rejection.
\end{itemize}

The Python integration script therefore plays a critical bridging role between the neural forecasting layer and the symbolic reasoning layer. It ensures that statistical anomaly outputs are converted into a symbolic representation that can be interpreted and constrained by the ASP rule base.

\section{ASP Rule Categories}

The ASP rule base is intentionally compact and organized around a small set of interpretable reasoning categories. These include:
\begin{enumerate}
	\item a \textbf{persistence filter}, which favors anomalies that remain active across multiple consecutive hours;
	\item \textbf{physics-based plausibility constraints}, which reject anomalies that violate simple domain assumptions;
	\item \textbf{co-anomaly market logic}, which allows selected price anomalies to be retained when accompanied by concurrent load anomalies;
	\item a \textbf{final refinement rule}, which combines persistence and veto conditions into the symbolic anomaly decision;
	\item \textbf{veto tracking rules}, which make rejected anomalies explainable through explicit reason labels.
\end{enumerate}

These rules do not attempt to represent the full complexity of the Austrian electricity market. Instead, they encode a targeted subset of persistence, physical plausibility, and cross-signal reasoning that improves anomaly interpretability and reduces implausible raw detections.

\section{Complete ASP Rule Base}

% Configuration for code wrapping and appearance
\lstset{
	language=Prolog,
	breaklines=true,                % Enables automatic line breaking
	breakatwhitespace=true,         % Only break lines at spaces if possible
	postbreak=\mbox{\textcolor{gray}{$\hookrightarrow$}\space}, % Adds a small arrow at the break
	basicstyle=\small\ttfamily,      % Slightly smaller font to fit more on a line
	columns=flexible,
	xleftmargin=2em,                % Indent the block for better visual structure
	frame=single,                   % Adds a frame to clearly define the page boundaries
	showstringspaces=false
}

\begin{lstlisting}[caption={market\_rules.lp (Corrected \& Enhanced for Explainable AI)}]
	% ---------------------------------------------
	% market_rules.lp (Corrected & Enhanced for Explainable AI)
	% ---------------------------------------------
	
	#const min_persist = 3.
	#const wind_ramp_limit = 500.
	
	% --- RULE 1: PERSISTENCE FILTER ---
	persistent(G, T) :- 
	anomaly(G, T), 
	anomaly(G, T-1), 
	anomaly(G, T-2).
	
	% --- RULE 2: PHYSICS-BASED CONSTRAINTS (Lowercase names) ---
	invalid_solar(T) :- 
	anomaly("cf_solar", T),
	pred(shortwave_radiation, R, T),
	R <= 0.
	
	excessive_wind_ramp(T) :-
	anomaly("cf_wind", T),
	pred(wind_mw, V1, T),
	pred(wind_mw, V0, T-1),
	|V1 - V0| > wind_ramp_limit.
	
	% --- RULE 3: CO-ANOMALY LOGIC (Price-Load Correlation) ---
	confirmed_economic_spike(T) :-
	anomaly("price", T),
	anomaly("load_mw", T).
	
	% --- RULE 4: FINAL REFINEMENT (The "Fusion") ---
	valid_anomaly(G, T) :-
	persistent(G, T),
	not holiday(T),
	not invalid_solar(T),
	not excessive_wind_ramp(T).
	
	% Priority Rule: Co-anomalies bypass persistence check
	valid_anomaly("price", T) :-
	confirmed_economic_spike(T),
	not holiday(T).
	
	% --- RULE 5: EXPLAINABLE AI (Veto Tracking) ---
	% Identify and categorize rejections for visualization
	vetoed(G, T) :- anomaly(G, T), not valid_anomaly(G, T).
	
	veto_reason(T, "night_solar") :- invalid_solar(T).
	veto_reason(T, "wind_ramp_limit") :- excessive_wind_ramp(T).
	veto_reason(T, "holiday_calendar") :- anomaly(G, T), holiday(T).
	
	% --- SHOW ATOMS ---
	#show valid_anomaly/2.
	#show vetoed/2.
	#show veto_reason/2.
\end{lstlisting}

\section{Interpretation of the Rule Base}

The ASP rule file reveals several important design choices in the proposed framework. First, anomaly acceptance is intentionally conservative: a raw anomaly must either satisfy persistence and avoid all veto conditions or, in the special case of price, receive support from concurrent load-anomaly evidence. Second, the rule base gives priority to simple and interpretable domain constraints rather than opaque post-hoc adjustments. Third, the presence of explicit veto-reason predicates makes the symbolic layer suitable for explainable anomaly analysis, since rejected anomalies can be associated with identifiable domain-based causes.

This appendix therefore provides the formal symbolic specification underlying the anomaly-refinement stage discussed in the main chapters of the thesis. It shows how persistence, physical plausibility, calendar context, and cross-signal support are combined in a compact reasoning layer that complements the neural forecasting backbone.